\documentclass[12pt,a4paper]{article}

\usepackage[utf8]{inputenc}
\usepackage[T1]{fontenc}
\usepackage{lmodern}
\usepackage{amsmath}
\usepackage{amssymb}
\usepackage{amsthm}
\usepackage{mathtools}
\usepackage[a4paper, top=22mm, bottom=22mm, left=22mm, right=22mm, headheight=14pt]{geometry}
\usepackage{enumitem}
\usepackage{booktabs}
\usepackage{xcolor}
\usepackage{fancyhdr}
\usepackage{titlesec}
\usepackage{microtype}
\usepackage{parskip}
\usepackage{mdframed}
\usepackage{tabularx}
\usepackage{array}
\usepackage[hidelinks]{hyperref}


\definecolor{primary}{RGB}{25, 80, 150}
\definecolor{secondary}{RGB}{90, 90, 90}
\definecolor{answergreen}{RGB}{230, 245, 230}
\definecolor{answerborder}{RGB}{60, 140, 60}
\definecolor{defblue}{RGB}{230, 238, 252}
\definecolor{defborder}{RGB}{50, 100, 200}
\definecolor{thmyellow}{RGB}{255, 248, 220}
\definecolor{thmborder}{RGB}{200, 160, 0}
\definecolor{noteorange}{RGB}{255, 243, 224}
\definecolor{noteborder}{RGB}{220, 120, 0}
\definecolor{formulabg}{RGB}{240, 240, 255}
\definecolor{cheatbg}{RGB}{248, 248, 248}


\titleformat{\section}
  {\large\bfseries\color{primary}}{}{0pt}{}[\titlerule]
\titlespacing{\section}{0pt}{16pt}{6pt}
\titleformat{\subsection}
  {\normalsize\bfseries\color{secondary}}{}{0pt}{}
\titlespacing{\subsection}{0pt}{10pt}{4pt}


%% Answer key box
\newmdenv[backgroundcolor=answergreen, linecolor=answerborder,
  linewidth=1.5pt, innerleftmargin=10pt, innerrightmargin=10pt,
  innertopmargin=6pt, innerbottommargin=6pt,
  skipabove=3pt, skipbelow=3pt]{answerbox}

%% Definition box
\newmdenv[backgroundcolor=defblue, linecolor=defborder,
  linewidth=1.5pt, innerleftmargin=10pt, innerrightmargin=10pt,
  innertopmargin=6pt, innerbottommargin=6pt,
  skipabove=3pt, skipbelow=3pt]{definitionbox}

%% Theorem box
\newmdenv[backgroundcolor=thmyellow, linecolor=thmborder,
  linewidth=1.5pt, innerleftmargin=10pt, innerrightmargin=10pt,
  innertopmargin=6pt, innerbottommargin=6pt,
  skipabove=3pt, skipbelow=3pt]{theorembox}

%% Note/warning box
\newmdenv[backgroundcolor=noteorange, linecolor=noteborder,
  linewidth=1.5pt, innerleftmargin=10pt, innerrightmargin=10pt,
  innertopmargin=6pt, innerbottommargin=6pt,
  skipabove=3pt, skipbelow=3pt]{notebox}

%% Formula box
\newmdenv[backgroundcolor=formulabg, linecolor=primary,
  linewidth=1pt, innerleftmargin=10pt, innerrightmargin=10pt,
  innertopmargin=6pt, innerbottommargin=6pt,
  skipabove=3pt, skipbelow=3pt]{formulabox}

%% MCQ option list
\newlist{mcqlist}{enumerate}{1}
\setlist[mcqlist]{label=\textbf{\Alph*.}, leftmargin=2.5em,
  topsep=2pt, itemsep=0pt, parsep=0pt}

%% Subpart list
\newlist{subpartlist}{enumerate}{1}
\setlist[subpartlist]{label=\textbf{\alph*)}, leftmargin=2.5em,
  topsep=4pt, itemsep=6pt}

%% Steps list
\newlist{stepslist}{enumerate}{1}
\setlist[stepslist]{label=Step \arabic*., leftmargin=3em,
  topsep=2pt, itemsep=4pt}

%% Display-style math inline
\everymath{\displaystyle}

\pagestyle{fancy}
\fancyhf{}
\fancyhead[L]{\small\itshape Functions — Self-Assessment Quiz}
\fancyhead[R]{\small\itshape Mathematics | Beginner Level}
\fancyfoot[C]{\small Page \thepage}
\renewcommand{\headrulewidth}{0.4pt}

\begin{document}
\begin{center}
{\LARGE\bfseries Functions}\\[6pt]
{\large Mathematics \textbar\ Beginner Level \textbar\ Self-Assessment Quiz}\\[4pt]
\end{center}
\noindent\rule{\linewidth}{0.6pt}\\[6pt]
\begin{center}
\textbf{Time:} 15 \qquad \textbf{Total Marks:} 57 \qquad \textbf{Questions:} 15
\end{center}
\vspace{6pt}
\noindent\textbf{How to use this quiz:}
\begin{enumerate}[leftmargin=2em, topsep=2pt, itemsep=0pt, parsep=0pt]
  \item Attempt all questions on your own before checking the answer key.
  \item Note any topics where you struggled and revisit them.
  \item Symbols and variables carry their standard meanings.
\end{enumerate}
\vspace{10pt}
\section*{QUESTIONS}
\noindent\rule{\linewidth}{0.4pt}
\vspace{4pt}
\subsection*{Section A: Multiple Choice Questions \normalfont\small (6 questions \(\times\) 3 marks = 18 marks)}
\vspace{4pt}
\noindent\textbf{1.} Which of the following relations represents a function? \hfill\textit{[3 marks]}
\begin{mcqlist}
  \item $\{(1,2), (2,3), (1,4)\}$
  \item $\{(1,2), (2,2), (3,2)\}$
  \item $\{(2,1), (2,2), (3,3)\}$
  \item $\{(x,y) \mid x^2 + y^2 = 1\}$
\end{mcqlist}
\vspace{6pt}
\noindent\textbf{2.} What is the domain of the function $f(x) = \sqrt{x-3}$? \hfill\textit{[3 marks]}
\begin{mcqlist}
  \item $x > 3$
  \item $x \ge 3$
  \item $x < 3$
  \item $x \le 3$
\end{mcqlist}
\vspace{6pt}
\noindent\textbf{3.} If $f(x) = 2x + 1$, what is $f(5)$? \hfill\textit{[3 marks]}
\begin{mcqlist}
  \item 6
  \item 7
  \item 10
  \item 11
\end{mcqlist}
\vspace{6pt}
\noindent\textbf{4.} Which of the following graphs represents a function? \hfill\textit{[3 marks]}
\begin{mcqlist}
  \item A vertical line
  \item A horizontal line
  \item A circle
  \item An ellipse
\end{mcqlist}
\vspace{6pt}
\noindent\textbf{5.} Given $f(x) = x^2$ and $g(x) = x+1$, what is $(f \circ g)(x)$? \hfill\textit{[3 marks]}
\begin{mcqlist}
  \item $x^2 + 1$
  \item $(x+1)^2$
  \item $x^2 + x + 1$
  \item $x(x+1)^2$
\end{mcqlist}
\vspace{6pt}
\noindent\textbf{6.} Which of the following functions is an odd function? \hfill\textit{[3 marks]}
\begin{mcqlist}
  \item $f(x) = x^2$
  \item $f(x) = \cos(x)$
  \item $f(x) = x^3$
  \item $f(x) = x^2 + x$
\end{mcqlist}
\vspace{6pt}
\subsection*{Section B: True or False Questions \normalfont\small (3 questions \(\times\) 2 marks = 6 marks)}
\vspace{4pt}
\noindent\textbf{1.} The range of the function $f(x) = x^2$ is all real numbers. \hfill\textit{[2 marks]}\\\noindent\quad\textit{(Circle:)} \quad True \quad / \quad False
\vspace{8pt}
\noindent\textbf{2.} The function $f(x) = 3x - 5$ is a linear function. \hfill\textit{[2 marks]}\\\noindent\quad\textit{(Circle:)} \quad True \quad / \quad False
\vspace{8pt}
\noindent\textbf{3.} If $f(x)$ is an even function, then its graph is symmetric with respect to the y-axis. \hfill\textit{[2 marks]}\\\noindent\quad\textit{(Circle:)} \quad True \quad / \quad False
\vspace{8pt}
\subsection*{Section C: Short Answer Questions \normalfont\small (3 questions \(\times\) 5 marks = 15 marks)}
\vspace{4pt}
\noindent\textbf{1.} Find the domain of the function $f(x) = \frac{1}{x-4}$. \hfill\textit{[5 marks]}
\vspace{22pt}
\noindent\rule{\linewidth}{0.2pt}
\vspace{4pt}
\noindent\textbf{2.} Given $f(x) = x^2 - 3x + 2$, evaluate $f(-1)$. \hfill\textit{[5 marks]}
\vspace{22pt}
\noindent\rule{\linewidth}{0.2pt}
\vspace{4pt}
\noindent\textbf{3.} Determine if the function $f(x) = x^4 - 2x^2$ is even, odd, or neither. \hfill\textit{[5 marks]}
\vspace{22pt}
\noindent\rule{\linewidth}{0.2pt}
\vspace{4pt}
\subsection*{Section D: Problem Solving / Applied Questions \normalfont\small (3 questions --- 18 marks total)}
\vspace{4pt}
\noindent\textbf{1.} A stone is thrown vertically upward from the ground with an initial velocity of $20 \text{ m/s}$. Its height $h$ (in meters) above the ground after $t$ seconds is given by the function $h(t) = 20t - 5t^2$.
\vspace{4pt}
\begin{subpartlist}
  \item Find the height of the stone after $1$ second. \hfill\textit{[2 marks]}
  \vspace{18pt}
  \noindent\rule{0.75\linewidth}{0.2pt}
  \item Find the height of the stone after $3$ seconds. \hfill\textit{[2 marks]}
  \vspace{18pt}
  \noindent\rule{0.75\linewidth}{0.2pt}
  \item At what time(s) is the stone at a height of $0$ meters (on the ground)? \hfill\textit{[2 marks]}
  \vspace{18pt}
  \noindent\rule{0.75\linewidth}{0.2pt}
\end{subpartlist}
\vspace{6pt}
\noindent\textbf{2.} Consider the functions $f(x) = 3x + 2$ and $g(x) = x^2 - 1$.
\vspace{4pt}
\begin{subpartlist}
  \item Calculate $f(g(2))$. \hfill\textit{[2 marks]}
  \vspace{18pt}
  \noindent\rule{0.75\linewidth}{0.2pt}
  \item Calculate $g(f(0))$. \hfill\textit{[2 marks]}
  \vspace{18pt}
  \noindent\rule{0.75\linewidth}{0.2pt}
  \item Find the expression for $(f+g)(x)$. \hfill\textit{[2 marks]}
  \vspace{18pt}
  \noindent\rule{0.75\linewidth}{0.2pt}
\end{subpartlist}
\vspace{6pt}
\noindent\textbf{3.} A small business produces custom-made T-shirts. The cost to produce $x$ T-shirts is given by the function $C(x) = 5x + 100$, and the revenue from selling $x$ T-shirts is given by $R(x) = 15x$.
\vspace{4pt}
\begin{subpartlist}
  \item What is the cost of producing $50$ T-shirts? \hfill\textit{[2 marks]}
  \vspace{18pt}
  \noindent\rule{0.75\linewidth}{0.2pt}
  \item Write the profit function $P(x)$. (Profit = Revenue - Cost) \hfill\textit{[2 marks]}
  \vspace{18pt}
  \noindent\rule{0.75\linewidth}{0.2pt}
  \item How many T-shirts must be sold for the business to break even (i.e., profit is $0$)? \hfill\textit{[2 marks]}
  \vspace{18pt}
  \noindent\rule{0.75\linewidth}{0.2pt}
\end{subpartlist}
\vspace{6pt}
\clearpage
\section*{ANSWER KEY}
\noindent\rule{\linewidth}{0.4pt}
\vspace{4pt}
\subsection*{Section A: Multiple Choice Questions \normalfont\small (6 questions --- 18 marks)}
\vspace{4pt}
\begin{answerbox}
\noindent\textbf{1.} \textbf{B}) $\{(1,2), (2,2), (3,2)\}$
\\ \noindent\textit{Explanation:} A relation is a function if each input (x-value) corresponds to exactly one output (y-value). Option A has '1' mapping to '2' and '4'. Option C has '2' mapping to '1' and '2'. Option D represents a circle, where for most x-values, there are two y-values. Option B satisfies the function definition as each x-value is unique or maps to a single y-value.
\end{answerbox}
\vspace{4pt}
\begin{answerbox}
\noindent\textbf{2.} \textbf{B}) $x \ge 3$
\\ \noindent\textit{Explanation:} For the square root function $f(x) = \sqrt{g(x)}$ to be defined in real numbers, the expression under the square root, $g(x)$, must be non-negative. Therefore, we must have $x-3 \ge 0$, which implies $x \ge 3$. This means the domain includes all real numbers greater than or equal to 3.
\end{answerbox}
\vspace{4pt}
\begin{answerbox}
\noindent\textbf{3.} \textbf{D}) 11
\\ \noindent\textit{Explanation:} To find $f(5)$, substitute $x=5$ into the function definition. So, $f(5) = 2(5) + 1$. This simplifies to $10 + 1 = 11$. Therefore, the value of the function at $x=5$ is 11.
\end{answerbox}
\vspace{4pt}
\begin{answerbox}
\noindent\textbf{4.} \textbf{B}) A horizontal line
\\ \noindent\textit{Explanation:} According to the vertical line test, a graph represents a function if and only if every vertical line intersects the graph at most once. A vertical line fails this test immediately. A circle and an ellipse also fail the test as many vertical lines intersect them at two points. A horizontal line, however, passes the vertical line test, representing a constant function.
\end{answerbox}
\vspace{4pt}
\begin{answerbox}
\noindent\textbf{5.} \textbf{B}) $(x+1)^2$
\\ \noindent\textit{Explanation:} The composition of functions $(f \circ g)(x)$ is defined as $f(g(x))$. First, substitute $g(x)$ into $f(x)$. Since $f(x) = x^2$ and $g(x) = x+1$, we replace the $x$ in $f(x)$ with $(x+1)$. Thus, $f(g(x)) = f(x+1) = (x+1)^2$.
\end{answerbox}
\vspace{4pt}
\begin{answerbox}
\noindent\textbf{6.} \textbf{C}) $f(x) = x^3$
\\ \noindent\textit{Explanation:} An odd function satisfies the property $f(-x) = -f(x)$. Let's test each option: For $f(x)=x^2$, $f(-x)=(-x)^2=x^2=f(x)$, so it's even. For $f(x)=\cos(x)$, $f(-x)=\cos(-x)=\cos(x)=f(x)$, so it's even. For $f(x)=x^3$, $f(-x)=(-x)^3=-x^3=-f(x)$, so it's odd. For $f(x)=x^2+x$, $f(-x)=(-x)^2+(-x)=x^2-x$, which is neither $f(x)$ nor $-f(x)$, so it's neither even nor odd.
\end{answerbox}
\vspace{4pt}
\subsection*{Section B: True or False Questions \normalfont\small (3 questions --- 6 marks)}
\vspace{4pt}
\begin{answerbox}
\noindent\textbf{1.} \textbf{False}
\\ \noindent\textit{Explanation:} The function $f(x) = x^2$ produces non-negative output values for any real input $x$. Since the square of any real number is always greater than or equal to zero, the range of $f(x)=x^2$ is $[0, \infty)$, not all real numbers. For example, $-1$ is not in the range.
\end{answerbox}
\vspace{4pt}
\begin{answerbox}
\noindent\textbf{2.} \textbf{True}
\\ \noindent\textit{Explanation:} A linear function is a polynomial function of degree one, which can be written in the form $f(x) = mx + b$, where $m$ and $b$ are constants. In this case, $f(x) = 3x - 5$ fits this form with $m=3$ and $b=-5$, confirming it is indeed a linear function.
\end{answerbox}
\vspace{4pt}
\begin{answerbox}
\noindent\textbf{3.} \textbf{True}
\\ \noindent\textit{Explanation:} An even function is defined by the property $f(-x) = f(x)$ for all $x$ in its domain. This means that if a point $(x, y)$ is on the graph, then the point $(-x, y)$ is also on the graph. This precisely describes symmetry with respect to the y-axis, where reflecting across the y-axis leaves the graph unchanged.
\end{answerbox}
\vspace{4pt}
\subsection*{Section C: Short Answer Questions \normalfont\small (3 questions --- 15 marks)}
\vspace{4pt}
\begin{answerbox}
\noindent\textbf{1.} The domain is $\{x \mid x \in \mathbb{R}, x \ne 4\}$ or in interval notation $(-\infty, 4) \cup (4, \infty)$.
\\ \noindent\textit{Explanation:} For a rational function, the denominator cannot be zero. In this case, the denominator is $x-4$. Setting $x-4 = 0$, we find that $x=4$. Therefore, the function is undefined when $x=4$. All other real numbers are valid inputs. So, the domain is all real numbers except 4. This can be expressed as $x \in \mathbb{R}, x \ne 4$ or using interval notation, $(-\infty, 4) \cup (4, \infty)$.
\end{answerbox}
\vspace{4pt}
\begin{answerbox}
\noindent\textbf{2.} $f(-1) = 6$
\\ \noindent\textit{Explanation:} To evaluate $f(-1)$, substitute $x = -1$ into the function's expression. So, $f(-1) = (-1)^2 - 3(-1) + 2$. First, calculate the square: $(-1)^2 = 1$. Next, calculate the product: $-3(-1) = 3$. Then, sum the terms: $1 + 3 + 2 = 6$. Therefore, $f(-1) = 6$.
\end{answerbox}
\vspace{4pt}
\begin{answerbox}
\noindent\textbf{3.} The function $f(x) = x^4 - 2x^2$ is an even function.
\\ \noindent\textit{Explanation:} To determine if a function is even, odd, or neither, we need to evaluate $f(-x)$. Substitute $-x$ for $x$ in the function: $f(-x) = (-x)^4 - 2(-x)^2$. This simplifies to $f(-x) = x^4 - 2x^2$. Since $f(-x)$ is equal to the original function $f(x)$, the function $f(x) = x^4 - 2x^2$ is an even function. This indicates symmetry about the y-axis.
\end{answerbox}
\vspace{4pt}
\subsection*{Section D: Problem Solving / Applied Questions \normalfont\small (3 questions --- 18 marks)}
\vspace{4pt}
\begin{answerbox}
\noindent\textbf{1.} See subparts
\\ \noindent\textit{Explanation:} Multi-part applied problem related to projectile motion described by a quadratic function.
\\\textbf{a)} $h(1) = 15 \text{ meters}$\\\textit{Explanation:} To find the height after $1$ second, we substitute $t=1$ into the function $h(t) = 20t - 5t^2$. So, $h(1) = 20(1) - 5(1)^2 = 20 - 5(1) = 20 - 5 = 15$. The height of the stone after 1 second is 15 meters.
\\\textbf{b)} $h(3) = 15 \text{ meters}$\\\textit{Explanation:} To find the height after $3$ seconds, substitute $t=3$ into the function $h(t) = 20t - 5t^2$. So, $h(3) = 20(3) - 5(3)^2 = 60 - 5(9) = 60 - 45 = 15$. The height of the stone after 3 seconds is 15 meters.
\\\textbf{c)} The stone is at $0$ meters height at $t=0 \text{ seconds}$ and $t=4 \text{ seconds}$.\\\textit{Explanation:} To find when the height is $0$ meters, set $h(t)=0$: $20t - 5t^2 = 0$. Factor out $5t$: $5t(4-t) = 0$. This gives two possible solutions: $5t=0 \implies t=0$ or $4-t=0 \implies t=4$. So, the stone is on the ground at the start ($t=0$) and after 4 seconds.
\end{answerbox}
\vspace{4pt}
\begin{answerbox}
\noindent\textbf{2.} See subparts
\\ \noindent\textit{Explanation:} This problem involves evaluating and composing different types of functions, specifically a linear function and a quadratic function.
\\\textbf{a)} $f(g(2)) = 11$\\\textit{Explanation:} First, evaluate $g(2)$. $g(2) = (2)^2 - 1 = 4 - 1 = 3$. Next, substitute this result into $f(x)$. So, $f(g(2)) = f(3) = 3(3) + 2 = 9 + 2 = 11$.
\\\textbf{b)} $g(f(0)) = 3$\\\textit{Explanation:} First, evaluate $f(0)$. $f(0) = 3(0) + 2 = 0 + 2 = 2$. Next, substitute this result into $g(x)$. So, $g(f(0)) = g(2) = (2)^2 - 1 = 4 - 1 = 3$.
\\\textbf{c)} $(f+g)(x) = x^2 + 3x + 1$\\\textit{Explanation:} The sum of two functions $(f+g)(x)$ is defined as $f(x) + g(x)$. Substituting the given functions, we get $(f+g)(x) = (3x+2) + (x^2-1)$. Combining like terms, the expression becomes $x^2 + 3x + (2-1) = x^2 + 3x + 1$.
\end{answerbox}
\vspace{4pt}
\begin{answerbox}
\noindent\textbf{3.} See subparts
\\ \noindent\textit{Explanation:} This problem involves applying linear cost and revenue functions to determine profit and break-even points for a business scenario.
\\\textbf{a)} The cost is $C(50) = 350 \text{ units of currency}$.\\\textit{Explanation:} To find the cost of producing $50$ T-shirts, substitute $x=50$ into the cost function $C(x) = 5x + 100$. So, $C(50) = 5(50) + 100 = 250 + 100 = 350$. The cost of producing 50 T-shirts is 350 units of currency.
\\\textbf{b)} $P(x) = 10x - 100$\\\textit{Explanation:} The profit function $P(x)$ is calculated as $P(x) = R(x) - C(x)$. Substitute the given revenue and cost functions: $P(x) = (15x) - (5x + 100)$. Distribute the negative sign: $P(x) = 15x - 5x - 100$. Combine like terms to get $P(x) = 10x - 100$.
\\\textbf{c)} The business must sell $10$ T-shirts to break even.\\\textit{Explanation:} To break even, the profit must be zero. Set the profit function $P(x) = 10x - 100$ equal to zero: $10x - 100 = 0$. Add $100$ to both sides: $10x = 100$. Divide by $10$: $x = 10$. Therefore, the business must sell 10 T-shirts to break even.
\end{answerbox}
\vspace{4pt}
\end{document}